% TODO checklist:
% - Inspected: Makefile (324 lines, all targets)
% - Inspected: scripts/deployment/deploy-production.sh
% - Inspected: scripts/auth/fetch_tokens.py
% - Inspected: scripts/ollama/init_default_model.py
% - Inspected: scripts/openapi/export_openapi.py (via directory listing)
% - Inspected: ops/backup/backup.sh
% - Inspected: scripts/database/ directory
% - Inspected: README.md operational scripts section

\section{Operational Scripts}
\label{sec:operational-scripts}

The Cineca Agentic Platform includes a comprehensive suite of operational scripts and Makefile targets that automate common development and deployment tasks. This section documents the available scripts and their usage.

\subsection{Makefile Targets}

The project Makefile provides the primary interface for common operations. Targets are organized by category:

\subsubsection{Environment and Setup}

\begin{lstlisting}[language=bash, caption=Environment setup targets]
# Create .env file from example template
make env

# Install Python dependencies
make install

# Install pre-commit hooks
make pre-commit-install
\end{lstlisting}

\subsubsection{Development Server}

\begin{lstlisting}[language=bash, caption=Development server targets]
# Run FastAPI with hot reload (local development)
make dev

# Probe health endpoints
make ready
\end{lstlisting}

\subsubsection{Docker Compose Operations}

\begin{lstlisting}[language=bash, caption=Docker Compose targets]
# Start all services
make up

# Start with CPU profile (extended timeouts)
make up-cpu

# Start with GPU profile (CUDA support)
make up-gpu

# Start only observability services
make up-observability

# Start only Redis service
make up-redis

# Stop all services
make down

# Stop and remove volumes (destructive)
make clean

# Restart app service
make restart

# View logs (all or specific service)
make logs         # All services
make logs S=app   # Specific service

# Show process status
make ps
\end{lstlisting}

\subsubsection{Code Quality}

\begin{lstlisting}[language=bash, caption=Code quality targets]
# Auto-format code (black + ruff)
make fmt

# Lint with auto-fixes
make lint

# Static type checking
make typecheck

# Run all quality checks
make check
\end{lstlisting}

\subsubsection{Testing}

\begin{lstlisting}[language=bash, caption=Testing targets]
# Run quick tests
make test

# Run full test suite (including e2e, performance)
make test-all

# Run Ollama-focused tests
make test-ollama

# Run Memgraph NL integration tests
make test-memgraph-nl

# Run smoke test (first 3 prompts)
make test-memgraph-nl-smoke

# Run security linters and scanners
make security

# Run CI tests in Docker
make test-ci
\end{lstlisting}

\subsubsection{Database Operations}

\begin{lstlisting}[language=bash, caption=Database operation targets]
# Memgraph operations
make seed      # Seed original DB
make populate  # Populate demo data
make backup    # Backup to backups/ directory
make restore BACKUP=backups/backup_xxx.tar.gz

# PostgreSQL operations
make db-migrate       # Run Alembic migrations
make db-migrate-down  # Rollback last migration
make db-reset         # Drop and recreate (destructive)
make db-seed          # Seed demo tenants
make db-revision MSG="description"  # Create new migration
make db-shell         # Open psql shell
make db-logs          # View PostgreSQL logs
\end{lstlisting}

\subsubsection{Model Operations}

\begin{lstlisting}[language=bash, caption=Model operation targets]
# Bootstrap model artifacts
make bootstrap-models TARGET_DIR=/opt/models

# Create Ollama models from Modelfiles
make ollama-models

# Verify LLM configuration
make llm-smoke-test

# Register providers and create instances
make runtime-smoke
\end{lstlisting}

\subsubsection{Documentation}

\begin{lstlisting}[language=bash, caption=Documentation targets]
# Export OpenAPI specification
make openapi

# Export OpenAPI via Docker container
make openapi-docker

# Aggregate CI target
make ci
\end{lstlisting}

\subsection{Deployment Scripts}

\subsubsection{Production Deployment Script}

The \texttt{scripts/deployment/deploy-production.sh} script automates production deployments:

\begin{lstlisting}[language=bash, caption=Production deployment script usage]
# Deploy to staging
./scripts/deployment/deploy-production.sh staging

# Deploy to production
./scripts/deployment/deploy-production.sh production

# Check current status
./scripts/deployment/deploy-production.sh status

# Rollback to previous version
./scripts/deployment/deploy-production.sh rollback
\end{lstlisting}

The script performs the following steps:
\begin{enumerate}
    \item Checks prerequisites (Docker, Docker Compose, \texttt{.env} file)
    \item Creates backup of current state
    \item Pulls latest images and rebuilds
    \item Deploys new containers
    \item Runs smoke tests to verify deployment
    \item Prompts for rollback if verification fails
\end{enumerate}

\subsubsection{Pre-flight Checks}

The \texttt{scripts/deployment/preflight\_checks.sh} script validates the deployment environment:

\begin{lstlisting}[language=bash, caption=Pre-flight checks script]
# Run pre-flight checks
bash scripts/deployment/preflight_checks.sh

# Checks performed:
# 1. Build & start services
# 2. Auth0 tokens script availability
# 3. CPU-only verification (no GPU dependencies)
# 4. MCP tools loaded (>=32 required)
# 5. Ollama models reachable
# 6. Intent keyword in test files
# 7. Script permissions
\end{lstlisting}

\subsection{Authentication Scripts}

\subsubsection{Token Fetching}

The \texttt{scripts/auth/fetch\_tokens.py} script retrieves Auth0 tokens for different roles:

\begin{lstlisting}[language=bash, caption=Token fetching script]
# Fetch all tokens (admin, user, machine)
python scripts/auth/fetch_tokens.py

# Tokens saved to:
# - /tmp/tokens.sh (shell export format)
# - /tmp/tokens.json (JSON format)

# Usage in tests
source /tmp/tokens.sh
curl -H "Authorization: Bearer $ADMIN_TOKEN" \
     http://localhost:8000/v1/health/ready
\end{lstlisting}

The script fetches three types of tokens:
\begin{itemize}
    \item \textbf{ADMIN\_TOKEN}: Full system access with \texttt{admin:all} scope
    \item \textbf{USER\_TOKEN}: Standard user access with \texttt{tools:invoke:basic} scope
    \item \textbf{MACHINE\_TOKEN}: Machine-to-machine access via client credentials
\end{itemize}

\subsection{Model Initialization}

The \texttt{scripts/ollama/init\_default\_model.py} script initializes the default LLM model:

\begin{lstlisting}[language=bash, caption=Model initialization script]
# Initialize default model in database
python scripts/ollama/init_default_model.py

# Environment variables:
# DEFAULT_MODEL_NAME: Model name (default: phi3:mini)
# OLLAMA_BASE_URL: Ollama API URL (default: http://ollama:11434/v1)
\end{lstlisting}

The script performs:
\begin{enumerate}
    \item Checks for existing healthy default model
    \item Creates Ollama provider if not exists
    \item Creates model instance if not exists
    \item Sets instance as global default
    \item Warms up the model with a test call
    \item Updates provider health status in Redis
\end{enumerate}

\subsection{Backup Scripts}

\subsubsection{Comprehensive Backup Script}

The \texttt{ops/backup/backup.sh} script provides automated backup for all databases:

\begin{lstlisting}[language=bash, caption=Backup script usage]
# Backup all databases
./ops/backup/backup.sh

# Backup specific database
./ops/backup/backup.sh --type postgres
./ops/backup/backup.sh --type redis
./ops/backup/backup.sh --type memgraph

# Backup and upload to S3
./ops/backup/backup.sh --type all --upload
\end{lstlisting}

Features include:
\begin{itemize}
    \item PostgreSQL backup with \texttt{pg\_dump}
    \item Redis backup via \texttt{BGSAVE}
    \item Memgraph snapshot creation
    \item Optional S3 upload for remote storage
    \item Retention policy enforcement (default: 30 days)
\end{itemize}

\subsubsection{Database-Specific Scripts}

\begin{lstlisting}[language=bash, caption=Database backup scripts]
# PostgreSQL backup
bash scripts/database/backup_database.sh

# PostgreSQL restore
bash scripts/database/restore_database.sh backup_file.sql.gz

# Validate PostgreSQL migration
bash scripts/database/validate_postgres_migration.sh
\end{lstlisting}

\subsection{OpenAPI Export}

The \texttt{scripts/openapi/export\_openapi.py} script exports the OpenAPI specification:

\begin{lstlisting}[language=bash, caption=OpenAPI export]
# Export via Python
python scripts/openapi/export_openapi.py

# Export via Makefile
make openapi

# Export via Docker (no local dependencies)
make openapi-docker

# Output files:
# - api/openapi.json (v1 endpoints)
# - api/openapi_v2.json (v2 endpoints)
\end{lstlisting}

\subsection{Developer Tooling and Test Modes}
\label{subsec:developer-tooling}

The platform provides several developer-focused utilities:

\subsubsection{Environment Doctor}

The \texttt{doctor} target displays environment information:

\begin{lstlisting}[language=bash, caption=Environment doctor target]
make doctor

# Output:
# Python: Python 3.11.x
# Pip:    pip 24.x
# Ruff:   0.x.x
# Black:  24.x.x
# Mypy:   1.x.x
# Bandit: 1.x.x
# DC:     Docker Compose version v2.x.x
\end{lstlisting}

\subsubsection{Debug Scripts}

The \texttt{scripts/debug/} directory contains debugging utilities:

\begin{lstlisting}[language=bash, caption=Debug scripts]
# Debug agent execution
python scripts/debug/debug_agent_run.py

# Debug JWT token
python scripts/debug/debug_token.py

# Test LLM configuration
python scripts/debug/test_llm_config.py

# Test Ollama directly
python scripts/debug/test_ollama_simple.py
\end{lstlisting}

\subsubsection{Verification Scripts}

The \texttt{scripts/verification/} directory contains verification utilities:

\begin{lstlisting}[language=bash, caption=Verification scripts]
# Verify default model configuration
bash scripts/verification/verify_default_model.sh

# Verify DMR (Default Model Resolver) system
python scripts/verification/verify_dmr_system.py

# Verify UI backend connectivity
bash scripts/verification/verify_ui_backend.sh
\end{lstlisting}

\subsubsection{NL-to-Cypher Test Mode}

For deterministic testing of the NL-to-Cypher pipeline:

\begin{lstlisting}[language=bash, caption=Enabling NL-to-Cypher test mode]
# Enable test mode with JSON hint file
export MEMGRAPH_NL_TEST_MODE=true
export MEMGRAPH_RESPONSE_MODE=fallback-only

# The test mode uses a JSON mapping file to return
# predetermined Cypher queries for given prompts,
# bypassing LLM generation for reproducible tests
\end{lstlisting}

\subsection{Scripts Summary Table}

Table~\ref{tab:scripts-summary} provides an overview of key operational scripts.

\begin{table}[ht]
\centering
\caption{Key operational scripts summary.}
\label{tab:scripts-summary}
\small
\begin{tabularx}{\textwidth}{lXl}
\hline
\textbf{Script} & \textbf{Purpose} & \textbf{Category} \\
\hline
deploy-production.sh & Automated production deployment with rollback & Deployment \\
preflight\_checks.sh & Pre-deployment environment validation & Deployment \\
production\_checklist.sh & Interactive production readiness checklist & Deployment \\
fetch\_tokens.py & Fetch Auth0 tokens for testing & Authentication \\
init\_default\_model.py & Initialize default LLM in database & Models \\
backup.sh & Multi-database backup with S3 support & Backup \\
export\_openapi.py & Export OpenAPI specification & Documentation \\
validate\_postgres\_migration.sh & Verify migration status & Database \\
\hline
\end{tabularx}
\end{table}

